% AulaTeX - maqueta inicial de construcción descendente
% Contrato futuro: el Agente puede investigar, redactar, evaluar y compilar a partir de este archivo.
\documentclass[12pt]{article}
\usepackage[utf8]{inputenc}
\usepackage[T1]{fontenc}
\usepackage[spanish]{babel}
\usepackage{enumitem}
\usepackage{hyperref}
\usepackage{longtable}

\title{derecho-a-la-seguridad-social-lde}
\author{AulaTeX}
\date{\today}

\begin{document}
\maketitle

\noindent\textbf{Nodo editorial:} derecho-a-la-seguridad-social-lde\\
\textbf{Padre editorial:} interinstitucional\\
\textbf{Modo de generación:} reforzar\\
\textbf{Destino:} UnADM\\
\textbf{Contrato futuro:} ready-for-agent

\section*{Ingesta base}
\begin{itemize}[leftmargin=*]
  \item Texto libre proporcionado: no
  \item Documento de apoyo proporcionado: no
\end{itemize}

\section*{Objetivo}
\begin{itemize}[leftmargin=*]
  \item Estructurar una base editorial para analizar el derecho a la seguridad social desde su marco constitucional, legal e institucional en Mexico.
  \item Completar la maqueta inicial de derecho-a-la-seguridad-social-lde.
  \item Establecer una base editorial para analizar el derecho a la seguridad social desde su dimensión constitucional, legal e institucional.
  \item Preparar una estructura compilable y reusable para actividad, reporte y presentación.
  \item Establecer una base editorial para analizar el derecho a la seguridad social desde su dimension constitucional, legal e institucional.
  \item Preparar una estructura compilable y reusable para actividad, reporte y presentacion.
\end{itemize}

\section*{Competencias}
\begin{itemize}[leftmargin=*]
  \item Identifica fundamentos juridicos del derecho a la seguridad social.
  \item Relaciona normativa aplicable con casos y problematica social concreta.
  \item Argumenta soluciones juridicas con soporte normativo y doctrinal.
  \item Identificar fundamentos jurídicos y alcances del derecho a la seguridad social.
  \item Relacionar normas, instituciones y casos con problemas sociales concretos.
  \item Argumentar propuestas con sustento normativo, jurisprudencial y doctrinal.
  \item Identificar fundamentos juridicos y alcances del derecho a la seguridad social.
\end{itemize}

\section*{Resultados esperados}
\begin{itemize}[leftmargin=*]
  \item Mapa normativo minimo de la materia elaborado y citado.
  \item Analisis escrito con tesis, argumentos y conclusiones verificables.
  \item Producto academico compilable en formato institucional UnADM.
  \item Esquema normativo mínimo validado y citado.
  \item Producto escrito con tesis, argumentos y cierre coherente.
  \item Entregables LaTeX compilables bajo estándar UnADM.
  \item Esquema normativo minimo validado y citado.
  \item Entregables LaTeX compilables bajo estandar UnADM.
\end{itemize}

\section*{Estructura sugerida}
\begin{itemize}[leftmargin=*]
  \item 1. Introduccion y delimitacion conceptual.
  \item 2. Fundamento constitucional del derecho a la seguridad social.
  \item 3. Desarrollo legal e institucional en Mexico.
  \item 4. Mecanismos de exigibilidad y desafios actuales.
  \item 5. Actividad integradora y criterios de entrega.
  \item 6. Conclusiones y referencias.
  \item 1) Introducción y delimitación conceptual.
  \item 2) Marco constitucional y convencional aplicable.
  \item 3) Desarrollo legal e institucional del sistema mexicano.
  \item 4) Exigibilidad, tutela y retos actuales.
  \item 5) Actividad integradora: consigna, entregable y rúbrica.
  \item 6) Conclusiones y referencias.
  \item 1) Introduccion y delimitacion conceptual.
  \item 5) Actividad integradora: consigna, entregable y rubrica.
\end{itemize}

\section*{Criterios de evaluación}
\begin{itemize}[leftmargin=*]
  \item Correccion conceptual y juridica.
  \item Sustento de fuentes y citacion consistente.
  \item Calidad argumentativa y estructura logica.
  \item Cumplimiento tecnico de formato y compilacion.
  \item Rigor conceptual y jurídico.
  \item Calidad y pertinencia de fuentes.
  \item Solidez argumentativa y claridad expositiva.
  \item Cumplimiento de formato editorial y técnico.
  \item Rigor conceptual y juridico.
  \item Cumplimiento de formato editorial y tecnico.
\end{itemize}

\section*{Bibliografía requerida}
\begin{itemize}[leftmargin=*]
  \item Constitucion Politica de los Estados Unidos Mexicanos.
  \item Ley del Seguro Social vigente.
  \item Bibliografia academica registrada en bibliografia-unadm.bib.
  \item Constitución Política de los Estados Unidos Mexicanos (vigente).
  \item Ley del Seguro Social (vigente).
  \item Fuentes académicas verificables en bibliografia-unadm.bib.
  \item Constitucion Politica de los Estados Unidos Mexicanos (vigente).
  \item Fuentes academicas verificables en bibliografia-unadm.bib.
\end{itemize}

\section*{Marcadores de investigación}
\begin{itemize}[leftmargin=*]
  \item unidad-1-concepto-y-naturaleza-juridica
  \item unidad-2-marco-constitucional-y-convencional
  \item unidad-3-regimenes-y-prestaciones
  \item unidad-4-justiciabilidad-y-casos
  \item fuentes-oficiales-prioritarias
  \item Definir fuentes, citas y vacíos de contenido antes de redactar.
  \item unidad-1-conceptos-fundantes-y-delimitaciones
  \item unidad-2-bloque-constitucional-y-convencional
  \item unidad-3-instituciones-regimenes-prestaciones
  \item unidad-4-justiciabilidad-criterios-y-casos
  \item control-vigencia-normativa-y-jurisprudencial
  \item depuracion-bibliografia-unadm-para-la-materia
\end{itemize}

\section*{Indicaciones editoriales por entregable}

\subsection*{Plantilla}
\begin{itemize}[leftmargin=*]
  \item \\textbackslash\\{\\}documentclass\\{article\\}
  \item \\% usar \\textbackslash\\{\\}input de plantilla institucional compartida
  \item \\textbackslash\\{\\}title\\{Derecho a la Seguridad Social (LDE)\\}
  \item \\textbackslash\\{\\}author\\{UnADM\\}
  \item \\textbackslash\\{\\}date\\{\\}
  \item Definir plantilla base con reglas editoriales, assets y referencias de estilo antes de redactar.
  \item \\% \\textbackslash\\{\\}input\\{plantilla-compartida-desde-base\\}
  \item \\% Estructura: portada, objetivo, competencias, resultados, desarrollo, cierre
  \item \\% \\textbackslash\\{\\}input\\{base/.../plantilla-compartida\\}
  \item \\% Estructura minima: portada, objetivo, competencias, resultados, desarrollo, cierre
  \item Prepara una raíz institucional reutilizable con carpetas por carrera, carpeta assets y referencias claras a plantillas LaTeX compartidas.
  \item Usar plantillas compartidas en base/ mediante input sin copiar archivos de plantilla.
  \item Respetar contrato latexmk: compilar pasando solo el .tex objetivo.
  \item Declarar bibliografia con archivo institucional y sin rutas absolutas locales.
  \item Organizar por ruta de materia dentro de UnADM con subcarpetas para evidencias y assets.
  \item Incluir archivos canonicos: reporte-*.tex, actividad-*.tex, presentacion-*.tex y COMPILACION.md.
  \item Separar contenido editorial (objetivos, criterios, consignas) de contenido investigado (fuentes, citas).
  \item Trabaja con la memoria fundacional consolidada como fuente base.
\end{itemize}

\subsection*{Actividad}
\begin{itemize}[leftmargin=*]
  \item \\% Secciones minimas: objetivo, instrucciones, producto, rubrica
  \item \\textbackslash\\{\\}section\\{Actividad integradora\\}
  \item \\textbackslash\\{\\}subsection\\{Objetivo\\}
  \item \\textbackslash\\{\\}subsection\\{Instrucciones\\}
  \item \\textbackslash\\{\\}subsection\\{Criterios de evaluacion\\}
  \item Desarrollar la actividad respetando objetivo, criterios de evaluación y marcadores de investigación.
  \item \\textbackslash\\{\\}subsection\\{Producto entregable\\}
  \item \\% Incluir rubrica con indicadores observables y ponderacion
  \item Estructurar una base editorial para analizar el derecho a la seguridad social desde su marco constitucional, legal e institucional en Mexico.
  \item Completar la maqueta inicial de derecho-a-la-seguridad-social-lde.
  \item Correccion conceptual y juridica.
  \item Sustento de fuentes y citacion consistente.
  \item Calidad argumentativa y estructura logica.
  \item Programa analitico oficial de la materia en UnADM.
  \item Marco constitucional mexicano del derecho a la seguridad social.
  \item Ley del Seguro Social y normativa complementaria vigente.
  \item Trabaja con la memoria fundacional consolidada como fuente base.
\end{itemize}

\subsection*{Reporte}
\begin{itemize}[leftmargin=*]
  \item \\% Secciones minimas de reporte juridico
  \item \\textbackslash\\{\\}section\\{Introduccion\\}
  \item \\textbackslash\\{\\}section\\{Marco normativo\\}
  \item \\textbackslash\\{\\}section\\{Analisis\\}
  \item \\textbackslash\\{\\}section\\{Conclusiones\\}
  \item \\textbackslash\\{\\}bibliography\\{bibliografia-unadm\\}
  \item Estructurar el reporte con bibliografía requerida y cierre verificable.
  \item \\textbackslash\\{\\}section\\{Marco conceptual y normativo\\}
  \item \\textbackslash\\{\\}section\\{Analisis juridico\\}
  \item \\% Verificar citas y correspondencia con .bib depurado
  \item \\% Verificar correspondencia completa entre citas y .bib
  \item Portada institucional y metadatos de asignatura.
  \item Objetivo, competencias y resultados esperados.
  \item Desarrollo tematico por unidades con subsecciones de analisis juridico.
  \item Redaccion formal, clara y orientada a aprendizaje juridico aplicado.
  \item Parrafos breves, terminologia legal precisa y definiciones operativas antes del analisis.
  \item Evitar afirmaciones dogmaticas sin sustento normativo o doctrinal verificable.
  \item Priorizar fuentes juridicas primarias: constitucion, leyes de seguridad social, reglamentos y criterios judiciales.
  \item Complementar con doctrina academica y textos de teoria juridica relevantes para UnADM.
  \item No inventar referencias; registrar solo fuentes verificadas en .bib.
  \item Toda afirmacion normativa debe vincularse a fuente juridica o academica verificable.
  \item La actividad debe incluir criterios de evaluacion observables y trazables a resultados esperados.
  \item Validar coherencia entre objetivo, competencias, producto entregable y rubrica.
\end{itemize}

\subsection*{Presentación}
\begin{itemize}[leftmargin=*]
  \item \\% Estructura breve para diapositivas
  \item \\textbackslash\\{\\}section\\{Contexto\\}
  \item \\textbackslash\\{\\}section\\{Fundamento juridico\\}
  \item \\textbackslash\\{\\}section\\{Hallazgos clave\\}
  \item \\textbackslash\\{\\}section\\{Cierre\\}
  \item Preparar la presentación con síntesis visual, continuidad editorial y evidencias clave.
  \item \\textbackslash\\{\\}section\\{Contexto del problema\\}
  \item \\textbackslash\\{\\}section\\{Hallazgos y analisis\\}
  \item \\textbackslash\\{\\}section\\{Conclusiones y recomendaciones\\}
  \item \\% Sintesis visual alineada al reporte y a la actividad
  \item Mapa normativo minimo de la materia elaborado y citado.
  \item Analisis escrito con tesis, argumentos y conclusiones verificables.
  \item Producto academico compilable en formato institucional UnADM.
  \item Redaccion formal, clara y orientada a aprendizaje juridico aplicado.
  \item Parrafos breves, terminologia legal precisa y definiciones operativas antes del analisis.
  \item Cubrir estructura base de actividad, reporte y presentacion sin desarrollar contenido final.
  \item Asegurar compatibilidad institucional UnADM y homogeneidad con nodos hermanos.
  \item Sintetiza visualmente la memoria editorial sin perder continuidad con plantilla, actividad y reporte.
\end{itemize}

\end{document}
